\section*{Capítulo \ref{ch:g11:functional-groups} --- Os grupos funcionais e as famílias}

\begin{solution}{exo:g11:functional-groups:1}
Hidroxila, \omterm{def:g11:functional-groups:alcohol}{álcool}; carbonila no meio da cadeia, \omterm{def:g11:functional-groups:carbonyl}{cetona}; carboxila,
\omterm{def:g11:functional-groups:carboxylic-acid}{ácido carboxílico}; grupo \omterm{def:g11:functional-groups:carboxylic-acid}{éster}, \omterm{def:g11:functional-groups:carboxylic-acid}{éster}; grupo amino \ce{-NH2}, \omterm{def:g11:functional-groups:amine}{amina}.
\end{solution}

\begin{solution}{exo:g11:functional-groups:2}
Metanol; propan-1-ol; ácido metanoico; propanal; propanona.
\end{solution}

\begin{solution}{exo:g11:functional-groups:3}
\omterm{def:g11:functional-groups:alcohol}{Álcool}; \omterm{def:g11:functional-groups:carbonyl}{aldeído}; \omterm{def:g11:functional-groups:carboxylic-acid}{éster}; \omterm{def:g11:functional-groups:amine}{amida}; \omterm{def:g11:functional-groups:halogenoalkane}{haloalcano}; \omterm{def:g11:functional-groups:halogenoalkane}{alceno}.
\end{solution}

\begin{solution}{exo:g11:functional-groups:4}
Propanal \ce{CH3-CH2-CHO}, propanona \ce{CH3-CO-CH3}. O propanal é o
\omterm{def:g11:functional-groups:carbonyl}{aldeído}: o carbono do seu \ce{C=O} está na ponta da cadeia, ligado a um
\omterm{def:g7:atoms-and-molecules:atom}{átomo} de hidrogênio; na propanona, ele está entre dois \omterm{def:g7:atoms-and-molecules:atom}{átomos} de
carbono.
\end{solution}

\begin{solution}{exo:g11:functional-groups:5}
Primário (o carbono que leva o \ce{-OH} tem um carbono vizinho);
secundário (dois); terciário (três).
\end{solution}

\begin{solution}{exo:g11:functional-groups:6}
\setchemfig{atom sep=1.5em}
Pentan-3-ol \chemfig{-[:30]-[:-30](-[:-90]OH)-[:30]-[:-30]};
ácido 2-metilbutanoico \chemfig{-[:30]-[:-30](-[:-90])-[:30](=[:90]O)-[:-30]OH};
hexan-2-ona \chemfig{-[:30](=[:90]O)-[:-30]-[:30]-[:-30]-[:30]};
etanoato de propila \chemfig{-[:30](=[:90]O)-[:-30]O-[:30]-[:-30]-[:30]}.
\end{solution}

\begin{solution}{exo:g11:functional-groups:7}
\ce{C3H6O}. Sim, eles são \omterm{def:g11:organic-skeletons:isomer}{isômeros}, mas com grupos diferentes: um
\omterm{def:g11:functional-groups:carbonyl}{aldeído} e uma \omterm{def:g11:functional-groups:carbonyl}{cetona} (os dois grupos carbonila, em lugares diferentes).
\end{solution}

\begin{solution}{exo:g11:functional-groups:8}
\ce{HCOO-CH2-CH3}, metanoato de etila; \ce{CH3-COO-CH2-CH2-CH3},
etanoato de propila.
\end{solution}

\begin{solution}{exo:g11:functional-groups:9}
Etanoato de metila (ácido etanoico, metanol); propanoato de etila (ácido
propanoico, etanol); metanoato de metila (ácido metanoico, metanol).
\end{solution}

\begin{solution}{exo:g11:functional-groups:10}
\setchemfig{atom sep=1.5em}
\begin{center}
\small
\begin{tabular}{cccc}
\chemfig{-[:30]-[:-30]-[:30]-[:-30]OH} &
\chemfig{-[:30](-[:90]OH)-[:-30]-[:30]} &
\chemfig{-[:30](-[:90])-[:-30]-[:30]OH} &
\chemfig{-[:0](-[:90]OH)(-[:-90])-[:0]} \\
butan-1-ol & butan-2-ol & 2-metilpropan-1-ol & 2-metilpropan-2-ol \\
primário & secundário & primário & terciário
\end{tabular}
\end{center}
\end{solution}

\begin{solution}{exo:g11:functional-groups:11}
Uma \omterm{def:g11:functional-groups:amine}{amida} (a etanamida). \emph{-eno}. Os \omterm{def:g11:functional-groups:carboxylic-acid}{ácidos carboxílicos} (\ce{C=O} e
\ce{-OH} no mesmo carbono), os \omterm{def:g11:functional-groups:carboxylic-acid}{ésteres} (\ce{C=O} e \ce{-O-}) e as \omterm{def:g11:functional-groups:amine}{amidas}
(\ce{C=O} e \ce{N}): três famílias.
\end{solution}

\begin{solution}{exo:g11:functional-groups:12}
Aspirina: um grupo carboxila (\omterm{def:g11:functional-groups:carboxylic-acid}{ácido carboxílico}) e um grupo \omterm{def:g11:functional-groups:carboxylic-acid}{éster}.
Paracetamol: um \ce{-OH} de fenol no anel e um grupo \omterm{def:g11:functional-groups:amine}{amida}
\ce{NH-CO}.
\end{solution}

\begin{solution}{exo:g11:functional-groups:13}
\ce{C2H6O}. Só o etanol, com o seu \ce{O-H}, forma \omterm{def:g11:polarity-and-cohesion:hydrogen-bond}{ligações de
hidrogênio} entre as suas \omterm{def:g7:atoms-and-molecules:molecule}{moléculas} (o metoximetano não tem hidrogênio no
seu oxigênio). O etanol ferve mais alto, a \qty{78}{\celsius}; o
metoximetano a \qty{-25}{\celsius}.
\end{solution}

\begin{solution}{exo:g11:functional-groups:14}
Um grupo hidroxila e um grupo carboxila. A cadeia tem três carbonos,
sendo o carbono do ácido o carbono 1 e o \ce{-OH} no carbono 2: ácido
2-hidroxipropanoico. Glicina: ácido 2-aminoetanoico (muitas vezes
escrito ácido aminoetanoico).
\end{solution}

\begin{solution}{exo:g11:functional-groups:15}
O etanol, \ce{CH3-CH2-OH} ou \ce{C2H6O}, é um \omterm{def:g11:functional-groups:alcohol}{álcool}; o etanal,
\ce{CH3-CHO} ou \ce{C2H4O}, um \omterm{def:g11:functional-groups:carbonyl}{aldeído}; o ácido etanoico,
\ce{CH3-COOH} ou \ce{C2H4O2}, um \omterm{def:g11:functional-groups:carboxylic-acid}{ácido carboxílico}. Do etanol ao etanal, perdem-se dois \omterm{def:g7:atoms-and-molecules:atom}{átomos} de hidrogênio;
do etanal ao ácido etanoico, ganha-se um \omterm{def:g7:atoms-and-molecules:atom}{átomo} de oxigênio.
\end{solution}

\begin{solution}{pb:g11:functional-groups:1}
\textbf{1.} A: grupo \omterm{def:g11:functional-groups:carboxylic-acid}{éster}; B: hidroxila; C: carboxila; D: carbonila na
ponta da cadeia; E: grupo \omterm{def:g11:functional-groups:carboxylic-acid}{éster}.

\textbf{2.} A \omterm{def:g11:functional-groups:carboxylic-acid}{éster}, B \omterm{def:g11:functional-groups:alcohol}{álcool}, C \omterm{def:g11:functional-groups:carboxylic-acid}{ácido carboxílico}, D \omterm{def:g11:functional-groups:carbonyl}{aldeído}, E \omterm{def:g11:functional-groups:carboxylic-acid}{éster}.

\textbf{3.} A (banana) e E (abacaxi): cheiros frutados.

\textbf{4.} Primário: o carbono que leva o \ce{-OH} está ligado a um
carbono.

\textbf{5.} \ce{C6H12O}; por exemplo a hexan-2-ona,
\ce{CH3-CO-CH2-CH2-CH2-CH3}.

\textbf{6.} A cadeia que contém o \ce{C-OH} tem quatro carbonos,
numerados a partir do \ce{-OH}, com um metil no carbono 3:
3-metilbutan-1-ol.

\textbf{7.} C: ácido etanoico; D: hexanal.

\textbf{8.} Butanoato de etila, aparentado ao ácido butanoico e ao
etanol.

\textbf{9.} \emph{Etanoato}: a parte ácida \ce{CH3-CO}, do ácido
etanoico; \emph{de 3-metilbutila}: a parte alcoólica, do
3-metilbutan-1-ol.

\textbf{10.} C (ácido etanoico) e B (3-metilbutan-1-ol).

\textbf{11.} A \ce{C7H14O2}; B \ce{C5H12O}; C \ce{C2H4O2}.

\textbf{12.} $\ce{C2H4O2} + \ce{C5H12O}$ contém \ce{C7H16O3}, uma
\ce{H2O} a mais que \ce{C7H14O2}: a água.

\textbf{13.} \ce{C2H4O2 + C5H12O -> C7H14O2 + H2O}: C 7 = 7, H 16 = 16,
O 3 = 3.

\textbf{14.} $M(\text{B}) = 5 \times 12.0 + 12 \times 1.0 + 16.0 =
\qty{88.0}{g/mol}$; $M(\text{C}) = \qty{60.0}{g/mol}$.

\textbf{15.} $60.0 + 88.0 - 18.0 = \qty{130.0}{g/mol}$.

\textbf{16.} $7 \times 12.0 + 14 \times 1.0 + 2 \times 16.0 =
\qty{130.0}{g/mol}$: a mesma.

\textbf{17.} E, \ce{C6H12O2}: $6 \times 12.0 + 12 \times 1.0 + 2 \times
16.0 = \qty{116.0}{g/mol}$.

\textbf{18.} \ce{C7H14O2}: a mesma fórmula que A, logo um \omterm{def:g11:organic-skeletons:isomer}{isômero}, com a
mesma \omterm{def:g10:the-mole:molar-mass}{massa molar}. Uma \omterm{def:g10:the-mole:molar-mass}{massa molar} sozinha não consegue distinguir
\omterm{def:g11:organic-skeletons:isomer}{isômeros}.

\textbf{19.} A e o ácido heptanoico têm os mesmos \omterm{def:g7:atoms-and-molecules:atom}{átomos}, mas grupos
diferentes, um \omterm{def:g11:functional-groups:carboxylic-acid}{éster} e um \omterm{def:g11:functional-groups:carboxylic-acid}{ácido carboxílico}: um tem cheiro de banana, o
outro de ranço. O grupo, e o lugar onde ele fica, decidem.

\textbf{20.} \qty{130.0}{g/mol}.
\end{solution}
